\documentclass[11pt]{article}
\usepackage[margin=27mm]{geometry}
\usepackage{amsmath,amssymb,amsthm,booktabs,hyperref}
\newtheorem{proposition}{Proposition}
\newcommand{\F}{\mathbb F}
\newcommand{\Sz}{\operatorname{Sz}}
\title{Two corrections to the diagrams for $\Sz(8)$}
\author{Human Compression Project}
\date{}
\begin{document}
\maketitle
\begin{abstract}
The subgroup descriptions in the rank-two Suzuki classification determine
the local residue orders without any graph computation. They also disprove
two quasi-primitivity labels under the definition printed in the source.
These corrections concern four numbered diagrams. The girths and diameters
of the twenty-nine diagrams require a separate finite verification.
\end{abstract}

\section{Residue orders and the printed orientation}
Let $G$ be a finite group, and let the two parabolics of a coset geometry be
$G_0,G_1$, with intersection $J$. Its type-$i$ elements are the left cosets
of $G_i$, so their number is $n_i=[G:G_i]$. The type-$i$ elements incident
with the fixed element $G_j$, for $j\ne i$, are precisely
\[
 \{gG_i:g\in G_j\}.
\]
Two representatives give the same coset precisely when they differ by an
element of $J$. The number of these elements is therefore $[G_j:J]$.
The source defines the order $s_i$ at its type-$i$ vertex as this residue
cardinality minus one [1, p.~100]. Thus
\begin{equation}
 (s_0,s_1)=([G_1:J]-1,[G_0:J]-1).                 \tag{1}
\end{equation}
The order at a diagram vertex is not the valency of that vertex in the
incidence graph: that valency counts elements of the other type.

\begin{proposition}
In source diagrams 25 and 26 [1, p.~118], the displayed orders $(12,4)$
must be replaced by $(4,12)$, with the printed parabolic orientation retained.
\end{proposition}
\begin{proof}
The left parabolic is $C_{13}\rtimes C_4$, of order $52$, and the right
parabolic is $C_5\rtimes C_4$, of order $20$. Their common subgroup has
order $4$. Formula (1) gives $20/4-1=4$ on the left and $52/4-1=12$
on the right. The printed element counts $(560,1456)$ independently fix
this orientation, since $|\Sz(8)|=29120$ and
\[
 29120/52=560,\qquad 29120/20=1456.
\]
The graph edge labels and the two parabolic labels are not changed by
this correction.
\end{proof}

\section{The quasi-primitivity labels}
The source calls a subgroup quasi-maximal when there is a unique chain
of intermediate subgroups between it and the ambient group, and calls
the geometry quasi-primitive when both parabolics have this property
[1, p.~99]. The following argument supplies three different saturated
chains for the parabolic in question.

Use the root coordinates already established in the normalizer treatment:
\[
 P=\{u(a,b):a,b\in\F_8\},\quad
 u(a,b)u(c,d)=u(a+c,b+d+ac^4).
\]
Let $Z=\{u(0,b):b\in\F_8\}$, and let $w=u(1,0)$.
Then $P$ has order $64$, $Z$ has order $8$, and
\[
 R=\{u(a,b):a\in\F_2,\ b\in\F_8\}
   \cong C_4\times C_2^2
\]
has order $16$. These identifications, and the normalizer
$N_G(\langle w\rangle)=R$, are acquired from the accepted root calculation.

\begin{proposition}
The $C_4\times C_2$ parabolic in source diagrams 27 and 29 is not
quasi-maximal in $\Sz(8)$. Both printed QPRI labels must be removed.
\end{proposition}
\begin{proof}
Write $H\cong C_4\times C_2$ for that parabolic. Its intersection with
the other parabolic is a cyclic subgroup $J$ of order four. All such
cyclic subgroups are conjugate, so conjugate the configuration until
$J=\langle w\rangle$. Since $H$ is abelian, it normalizes $J$, and hence
$H\le R$. Its order is $8$, so $[R:H]=2$.

The first-coordinate map $P\to(\F_8,+)$ is a surjective homomorphism
with kernel $Z$. The subgroup $R$ is its inverse image of $\F_2$.
Consequently $R\lhd P$ and
\[
 P/R\cong (\F_8/\F_2,+)\cong C_2^2.
\]
The latter group has three different subgroups of order two. Their inverse
images are three different subgroups $P_1,P_2,P_3$ of order $32$.
For the point stabilizer $B=P\rtimes C_7$, we obtain the three chains
\[
 H<R<P_i<P<B<G,\qquad i=1,2,3.
\]
The first three inclusions each have index two, $[B:P]=7$, and $B$ is a
maximal subgroup of $G$ by the acquired Suzuki subgroup theorem. Thus each
chain is saturated and the chains are distinct. This disproves uniqueness
even if ``chain'' is read as ``maximal chain''. One parabolic failing the
definition suffices to remove QPRI from the geometry.
\end{proof}

The proof uses the source's own definition and the already acquired root
and maximal-subgroup results. It does not require a subgroup census. The
local double-transitivity label on diagram 29 is independent of this
correction and is retained by the general classification theorem.

\begin{thebibliography}{9}
\bibitem{source} D. Leemans, \emph{The rank 2 geometries of the simple Suzuki
groups $\Sz(q)$}, Beitr\"age zur Algebra und Geometrie \textbf{39}
(1998), 97--120.
\end{thebibliography}
\end{document}
